% Transcription — fonds Alexandre Grothendieck, shelfmark 156-4, batch 1 (pages 1-20).
% Established from the facsimile archives/batches/156-4/batch-01.pdf (a copy of
% 156-4.pdf fetched through the project's relay; no cover sheet in the batch, so
% PDF page k = folder page k), per the skill transcribe-grothendieck. The folder
% is eighty-eight pages in five batches; this is the first.
%
% Pass: Opus 5.5 (claude-opus-5-5), 2026-09-26 — first pass, unchecked against the pages by a human.
%
% Pages skipped: none. Pages summarised: none. All twenty leaves are
% mathematics in his hand and all are transcribed.
%
% His pagination 1-20 runs across the top of the leaves and coincides with the
% archivists'. Page 1 carries the heading « Analysis situs (première mouture) »,
% under it « (10 Juin 1986) », and top right « GF IV » (Géométrie des formes,
% chapter IV). A second date, « (11 juin) », underlined, stands in the left
% margin of page 18. Both are given in \note{}s where they occur. His numbered
% runs: « 1. » and « I) Figures » on page 1 (conditions a)-d), f), and a
% marginal e)); the structure data (1)-(4) on page 9 (circled), extended by (5)
% on page 15 and a cancelled (6); axioms Fig_1 (p. 10), Fig_2 (p. 11); for a
% « contrée », data 1°)-4°) (p. 17) and axioms Contr_1 (p. 17), Contr_2
% (p. 18, a first version cancelled), Contr_3 (p. 20, begun and cancelled).
% The bracketed « [Chapitre] » of the inventory title is the archivists'.
%
% What the batch is about. The first draft of chapter IV. The model is a
% (tame) topological space: a « figure » is a closed part X with a locally
% finite decomposition into closed strata X_i (non-empty, locally connected,
% X_i rare in X_j for X_i ⊊ X_j, ∂X_i the union of smaller strata, int X_i a
% connected manifold, X_i ∩ X_j a union of strata), identified with its set
% of strata, so 𝔉 ⊂ 𝔓(𝔓_f(X)); Σ_F the ordered set of strata and Σ = ∐ Σ_F
% (p. 2). Sub-figures correspond to down-closed parts of Σ_F; « figures
% élémentaires » (one greatest stratum) form an ordered set 𝔉_él, and a figure
% is recovered as the closed part Σ'_F of 𝔉_él (pp. 3-4). F-saturated parts
% and compatibility of two figures (open strata disjoint or equal, a) ⇔ b)
% sketched by induction on dimension; pp. 5-7), configurations, locally finite
% configurations as an antifilter, and the resulting axiomatic structure (1)-(4)
% on 𝔉_él (pp. 7-10); axioms Fig_1, Fig_2 and combinatorial dimension,
% dimension-0 figures = submanifolds (pp. 10-12). Subdivision and a second
% order « F' ≪ F » (F' raffine F) on figures, with subdivision = maximal
% refinement of same support, and an attempted justification that runs into
% « un pb, qui a l'air intéressant » of topology of a stratified space inside a
% manifold (pp. 13-16). Page 17 abstracts a « contrée » from (𝔉_él, ≪, 𝓡,
% locally finite parts); geometric dimension via chains X_0 < … < X_{n-1} ≪ X_n
% (p. 18); page 19 renames « figure élémentaire » as « multistrate » and
% defines the « lieux » as the ≪-minimal multistrates, ℒ ⊂ 𝓜; page 20 sets
% up ℒ_X, ℒ̂_X and ℒ_F, ℒ̂_F to reconstruct ≪ and compatibility from them, and
% breaks off on a cancelled Contr_3.
%
% Notation fixed once. The ambient space of page 1, a boxed/crossed X, is set
% \mathbb{X}; from page 10 on the elementary figures (multistrates) are a
% barred X, set \mathbf{X} (the bar is not always written, and plain X is kept
% where he writes plain X). 𝔉, 𝔉_él, 𝔓 are \mathfrak; his barred cursive
% capital for a part of 𝔉_él is \Phi; compatibility is \mathcal{R}; the
% lieux and multistrates \mathcal{L}, \mathcal{M}; the figure on page 20
% \mathcal{F}. His underlining is set \emph. His circled numerals are given as
% (n) with a note.
%
% Hard pages: 5 (top half cancelled, long unread margin), 6 (the proof of
% a) ⇒ b), heavily corrected), 16 (full-margin note unread), 18 (lower half a
% palimpsest of oblique insertions, read only as a main thread). Readings
% flagged and not settled include « contrée » vs its first use on p. 17 (settled
% by p. 1, where it is clear), « décomposées » (p. 12), « se déduit » (p. 14),
% « canonique » (p. 16). Variances on the page kept as written: « {X_i} ∈
% 𝔓(Y) » for a part called X (p. 1); « raffine F' » for raffine F (p. 15);
% 𝔉 for 𝔉_él (p. 12); 𝔓(𝔓(X)) for 𝔓(𝔓(ℒ)) (p. 20).

\documentclass[11pt,a4paper]{article}
\input{../preamble/grothendieck}

\folder{156-4}
\batch{1}
\pages{1}{20}
\dating{1986}
\watermark{Édition de démonstration}
\foldertitle{[Chapitre] IV. Analysis situs (première mouture) : notes manuscrites (10/06/1986)}

\begin{document}
\page{1}
\section{Analysis situs (première mouture)}

\note{titre de sa main, souligné, en tête de la page 1, suivi de « (première mouture) » ; dessous, la date de sa main « (10 Juin 1986) ». Dans l'angle supérieur droit, dans un angle tracé à la plume, « GF IV » (le chiffre souligné) : Géométrie des formes, chapitre IV. Il numérote ses pages de 1 à 20 en haut des feuillets ; sa numérotation coïncide ici avec celle des archivistes}

1. Comme modèle « ensembliste topologique » de l'axiomatique d'un « modèle de formes \add{topologiques} » (cas d'une « \emph{contrée} »), considérons les notions suivantes. Modèle : un « espace » $\pm$ ordinaire.
\note{« topologiques » est écrit sous la ligne et renvoyé par un trait après « formes »}

On suppose ici que $\mathbb{X}$ \supplied{est} un espace topologique séparé, \uncertain{ou} un espace modéré (pas néc. compact) pour une théorie modérée donnée. Dans ce cas « partie fermée » signifie « partie fermée \uncertain{modérée} ».
\note{l'espace ambiant est noté par un X double, encadré ou barré, rendu ici $\mathbb{X}$ ; le mot sous « partie fermée », à la fin de la ligne, est surchargé}

I) \emph{Figures}. Une \emph{figure} $F$ / de $\mathbb{X}$ est une partie fermée $X$, \emph{munie} d'une décomposition de $Y$ en \struck{\ill{}} \add{parties} $X_i$ :
\[
  \{X_i\} \in \mathfrak{P}(Y) \qquad X = \bigcup X_i
\]
appelées les \emph{strates fermées}. On fait les hypothèses
\note{« de $Y$ », « $\mathfrak{P}(Y)$ » : sic, pour la partie fermée appelée $X$ une ligne plus haut}
\marginal{$X$ est appelé \struck{\ill{}} le \uncertain{support} de la figure $|F|$}
\note{note oblique dans la marge gauche, en face de la définition ; lue en partie}

\begin{itemize}
  \item[a)] Les $X_i$ fermés \add{non vides et loc. connexes} (pas néc. \uncertain{connexes}) \struck{\ill{}} \ill{}
  \marginal{(loc. c'est automatique si $\mathbb{X}$ l'est, \ill{} i.e. $\mathbb{X}$ \ill{})}
  \note{la note de marge est entre accolades sur la page}
  \item[b)] La famille des $X_i$ est loc. finie
  \item[c)] Si \struck{\ill{}} \add{(dans \uncertain{hyp.} a) b))} \struck{\ill{}} \add{$\Rightarrow$ réunion \ill{} de strates est fermée} $X_i \subset X_j$, $X_i \neq X_j$ (i.e. $i < j$) alors $X_i$ \emph{rare} dans $X_j$
  \item[d)] Posant $\partial X_i = \struck{\ill{}} \bigcup_{j < i} X_j$ (dans $X_i$
\end{itemize}
\note{les ajouts à a) et c) sont écrits au-dessus de la ligne, serrés contre le bord droit ; la fin de la ligne a) est surchargée et n'est lue qu'en partie}
\marginal{NB on pose $i \leq j \Leftrightarrow X_i \subset X_j$, \quad $i < j \Leftrightarrow X_i \subsetneq X_j$}

\page{2}
est fermé dans tous \uncertain{les} $X_i$, par a) b) c)) alors
\[
  \mathrm{int}(X_i) = X_i \smallsetminus \partial X_i
\]
(qui est donc un ouvert \add{rel.} dans $X_i$ \add{\ill{} $(X_j)$}, et non vide puisque $X_i$ l'est) \struck{\ill{}} \add{soit} 1°) une \emph{variété} 2°) \emph{connexe}. (NB dans le cas $\mathbb{X}$ modéré, on prend « variété » au sens modéré, \uncertain{entendu}.)
\note{« soit » et « $(X_j)$ » sont écrits au-dessus de la ligne ; « variété » est souligné deux fois}

f) $X_i \cap X_j$ est réunion \add{de} $F$, $X_k$.
\note{lettre « f) » écrite sur un autre signe ; le « $F$ » gras est écrit au-dessus de « de », sans doute pour « de strates de $F$ »}
\marginal{e) Si \struck{$X_i < X_j$}, alors \uncertain{dim} \ill{} $(X_{ij})$ \ill{} \uncertain{les} \ill{}}
\note{la condition e) est écrite en oblique dans la marge gauche, en partie biffée, avec un trait fléché et les signes « $+1$ », « $+0$ » près d'un gribouillis ; elle n'est lue que par fragments}

Donc une figure \add{$F$} peut s'identifier à l'ens. de ses strates, qui est un ens. de parties fermées de $X$
\[
  F \subset \mathfrak{P}_f(X) \qquad \text{i.e.} \qquad F \in \mathfrak{P}(\mathfrak{P}_f(X))
\]
($X$ se reconstituant \add{$\{X_i\}$} comme $\bigcup X_i$), satisfaisant certaines conditions \add{(qui \uncertain{vérifient} \ill{}, a) à d))}. On a donc l'ens. des figures, \struck{\ill{}} soit $\mathfrak{F}$, qui est donc un ens. de parties de $\mathfrak{P}_f(X)$
\[
  \mathfrak{F} \subset \mathfrak{P}(\mathfrak{P}_f) \qquad \text{i.e.} \qquad \mathfrak{F} \in \mathfrak{P}(\mathfrak{P}(\mathfrak{P}_f(X)))
\]
\note{$\mathfrak{P}_f(X)$ : l'indice $f$ semble désigner ici les parties fermées ; à la page 2 le $X$ de $\mathfrak{P}_f(X)$ est écrit en gras}

À chaque $F \in \mathfrak{F}$ est associé l'ens. \emph{ordonné} $\Sigma_F$ de ses strates. L'ens. somme
\[
  \Sigma = \coprod_{F \in \mathfrak{F}} \Sigma_F \longrightarrow \mathfrak{F}
\]
s'identifie à l'ens. des couples $(F, X_i)$, \struck{\ill{}} où $F$ est une figure et $X_i$ est une strate de $F$ i.e. $X_i \in F$.
\note{la flèche est verticale sur la page, de $\Sigma$ vers $\mathfrak{F}$}
\marginal{les fibres sont les $\Sigma_F$}
\marginal{\emph{NB} Dans le cas d'espèce, \ill{} à $\Sigma_F = F$ \struck{mais \ill{} \ill{}}}

\page{3}
Si $F$, $F'$ sont deux figures, on dit que $F'$ est une \emph{sous-figure} de $F$, si toute strate de $F'$ est une strate de $F$, i.e. $F' \subset F$. Donc \struck{le support} \uncertain{on a} \struck{$\sigma(F)$} et
\[
  |F'| \subset |F|
\]
\emph{NB} Si $X_i \in F'$, et $Y_j \in F$ est tel que $Y_j \subset X_i$, alors $Y_j \in F'$ \struck{En effet \ill{} $Y \in$ \ill{} $(Y_j)$, alors} \struck{La relation}

On écrit $F' \leq F$ pour cette relation, qui est une relation d'ordre dans $\mathfrak{F}$. L'ens. des sous-figures de $F$ est en corr. 1-1 avec l'ens. des parties $\Phi$ \add{($\subset \Sigma_F$)} de \struck{$\Sigma$} $\Sigma_F$ \add{(ens. des strates)} qui ont la propriété que \add{($i, j \in \Sigma_F$, $i \in \Phi$)}
\[
  j \leq i \Longrightarrow j \in \Phi
\]
(idéaux \struck{\ill{}} « \uncertain{fermés} » de l'ens. \struck{\ill{}} ordonné, pour la top. \ill{} \add{sur $\Sigma_F = \Sigma$} \struck{de fermés} \struck{d'ouverts} / \add{décroissants les} les $\Sigma_{\leq i}$).
\note{$\Phi$ rend une capitale cursive barrée, la même qu'aux pages 5, 8-11 pour une partie de $\mathfrak{F}_{\mathrm{él}}$. La fin de la parenthèse est surchargée : l'indice de $\Sigma_{\leq i}$ est écrit sur un signe biffé}

\struck{Ainsi, on a une relation d'ordre sur}
\marginal{\uncertain{11.6.} \uncertain{plus loin} \quad on dira « \uncertain{multistrate} »}
\note{dans la marge gauche, en oblique, « 11.6. » souligné, puis quelques mots ; ils semblent renvoyer au 11 juin (page 18) et au terme « multistrate » que la page 19 substitue à « figure élémentaire »}
\emph{Figure élémentaire} : qui a une strate maximale i.e. $\Sigma_F$ a un plus grand élément. Soit $\mathfrak{F}_{\mathrm{él}}$ l'ens. des figures élémentaires. \struck{Alors}
\note{$\mathfrak{F}_{\mathrm{él}}$ est entouré d'un cercle}

Pour \uncertain{toute} figure $F$, les sous-figures élémentaires sont en corr. 1-1 avec les

\page{4}
strates de $F$. \struck{Donc on peut dire}
\begin{itemize}
  \item[] $F$ figure
  \item[] $\Sigma_F$ ens. de ses strates
  \item[] $\Sigma'_F$ ens. des sous-figures élémentaires.
\end{itemize}
On a :
\[
  \Sigma'_F \simeq \Sigma_F, \qquad \Sigma'_F \subset \mathfrak{F}_{\mathrm{él}}
\]
\note{l'inclusion $\Sigma'_F \subset \mathfrak{F}_{\mathrm{él}}$ est écrite verticalement sous $\Sigma'_F$}
$F$ est connu quand on connaît $\Sigma'_F \subset \mathfrak{F}_{\mathrm{él}}$, car $F$ est l'ens. des strates des $F' \in \Sigma'_F$.

Mais la relation d'inclusion entre figures s'interprète comme une relation d'inclusion entre parties de $\mathfrak{F}_{\mathrm{él}}$.

\struck{Donc} Considérons la relation d'inclusion $\leq$ entre figures élémentaires. \struck{Alors} \add{Ainsi,} $\mathfrak{F}_{\mathrm{él}}$ est un ens. ordonné. \add{Si} \struck{Donc} $F$ est une figure, $\Sigma'_F$ est une \struck{figure dont \uncertain{dites}} partie fermée de l'ens. ordonné $\mathfrak{F}_{\mathrm{él}}$. L'inclusion entre figures est l'inclusion entre ces parties fermées.
\note{au-dessus de la ligne biffée, « $F$ est une figure » est récrit ; « Ainsi » et « Si » sont écrits au-dessus des mots biffés}

\page{5}
\struck{Pour qu'une partie $\Phi$ de $\mathfrak{F}_{\mathrm{él}}$ soit de la forme $\Sigma'_F$, il faut et \uncertain{suffit} :}

\struck{a) Que $\Phi$ soit une partie fermée de $\mathfrak{F}_{\mathrm{él}}$}

\struck{b) Que la \struck{famille} \ill{} \ill{} $|\varphi|$, \ill{}}

\struck{\add{$F$, $F'$} Deux figures \ill{} \uncertain{dites}}

\struck{Considérons deux figures $F$, $F'$, \uncertain{vérifiant} les conditions $|F| \cap |F'|$}
\note{le haut de la page, jusqu'à « $|F| \cap |F'|$ », est barré de longs traits obliques, les dernières lignes étant de plus enfermées dans un contour}

Soit $F$ une figure, une partie $A$ de $|F| = \mathrm{supp}\, F$ est dite \struck{\ill{}} \emph{$F$-saturée}, si saturée (pour $F$) i.e. elle est saturée pour la relation d'équivalence définie par les \struck{$X_i$} $\mathrm{int}(X_i)$ \add{(strates « ouvertes »)}. Si \ill{} $(X_i)$ est l'ens. des strates), cela signifie que \ill{} réunion de strates (fermées)

\marginal{\ill{} relation \ill{} $\cap$, $\cup$ \ill{} $|F|$ \ill{} $A$ \ill{} $F'$ \ill{} $A$ \ill{} $|F|$ \ill{} figure \ill{} $F' \subset (F')$ \ill{} $|F'|$ \ill{}}
\note{une longue note oblique couvre la marge gauche du bas de la page ; seuls des fragments en sont lus}

Soient $F$, $F'$ deux figures. Considérons les conditions

\struck{a) $|F| \cap |F'|$ est $F$-saturée et $F'$-saturée}

a) $\forall\, X \in \Sigma_F$, $X' \in \Sigma_{F'}$, $X \cap X'$ est \struck{\ill{}} $F_X$-saturée \struck{dans $F_X$} et $F'_{X'}$-saturée

\page{6}
\struck{c) $\exists$ \struck{sous-}figure $F''$, telle que}

\struck{1°) $F''$ sous-figure de $F$, et de $F'$}

\struck{2°) \struck{\ill{}} $(|F| \cap |F'|) = |F''|$}

\struck{Si \struck{$X$} $U$ strate \uncertain{connexe} de $F$, $U'$ strate}
\note{ces lignes sont barrées de traits obliques}

b) Si $U_i$ est une \struck{cellule} \add{strate} ouverte de $F$, $U'_j$ strate ouverte de $F'$, alors on a
\[
  U_i \cap U'_j = \emptyset \quad \text{\emph{ou}} \quad U_i = U'_j .
\]

Ces conditions sont équivalentes

\struck{b $\Rightarrow$ a trivial (puisque $|F| = \bigcup X_i$, $|F'| = \bigcup X'_j$)}

\struck{\ill{}} b) $\Longrightarrow$ a) trivial, prouvons

a) $\Longrightarrow$ b) \struck{Si} Supposons que \add{prouvons $U_i = U'_j$,} $U_i \cap U'_j \neq \emptyset$, \struck{$U_i \neq U'_j$ donc} \struck{\ill{} $X_i \neq X'_j$ Considérons $X_i \cap X'_j$ ($\neq \emptyset$),} soit \add{$x \in$} $U_i \cap U'_j$, Considérons $X_i \cap X'_j$, est \add{fermée} $F_{X_i}$-saturée, \struck{\ill{}} \add{et} \struck{et le \ill{}} \ill{} $x \in U_i$, \ill{} $X_i \subset X_i \cap X'_j$, \struck{\ill{} $X_i \cap X'_j$} i.e. $X_i \subset X'_j$. \ill{} $= X'_j \subset X_i$
\[
  X_i = X'_j .
\]
\struck{Donc} On en conclut \ill{} directement \struck{\ill{}} $\partial X_i = \partial X'_j$. Mais \struck{$X_i$} $U_i = U'_j$, \struck{on a} \add{i.e. a) $\Rightarrow$ b)} On \struck{le} prouve \add{par récurrence sur} $\sup(\dim F, \dim F')$ (on est \struck{\ill{}} \struck{ramené} \add{ramené} au cas des figures $F_{X_i}$, $F'_{X'_j}$ \add{de dim.} i.e. \uncertain{finie}\ldots)
\note{le raisonnement est écrit vite et beaucoup corrigé : plusieurs débuts de ligne, dans la marge gauche, sont biffés ; l'ordre des ajouts interlinéaires n'est pas sûr}

\page{7}
On dit que $F$, $F'$ sont \emph{compatibles} si les conditions en question sont satisfaites. Pour ceci, il faut et il suffit que pour toutes $F_0 \in \Sigma'_F$, $F'_0 \in \Sigma'_{F'}$, $F_0$ et $F'_0$ soient compatibles.

Donc, dans l'ens. \add{ordonné $\mathfrak{F}_{\mathrm{él}}$} des figures élémentaires, on a introduit une relation symétrique, réflexive)
\[
  \mathcal{R} \subset \mathfrak{F}_{\mathrm{él}} \times \mathfrak{F}_{\mathrm{él}}
\]
\marginal{\struck{\ill{}} $F_0$ comp. avec $F'_0$}
ou encore
\[
  S \subset \mathfrak{P}_2(\mathfrak{F}_{\mathrm{él}})
\]
\marginal{paires compatibles}
\note{les deux gloses sont écrites à droite des formules ; devant la première, un mot est raturé}

\emph{NB} Si \struck{\ill{}} $F_0$, $F'_0$ sont des fig. él. telles que $F_0 \leq F'_0$, alors elles sont compatibles. Mieux, si $F'_0, F''_0 \leq F$, \uncertain{alors} $F_0$ et $F'_0$ sont compatibles entre elles
\note{sic : « $F_0$ et $F'_0$ » après « si $F'_0, F''_0 \leq F$ » ; la fin de la phrase est écrite sous la ligne}

\emph{Configuration} Ens. de figures deux à deux compatibles. Configuration élémentaire : ens. de figures élémentaires deux à deux compatibles.

\page{8}
Soit $F$ une figure, \struck{$\Sigma'$}
\[
  \Sigma'_F = \Phi \subset \mathfrak{F}_{\mathrm{él}}
\]
la partie de $\mathfrak{F}_{\mathrm{él}}$ associée. \struck{Alors} On a les conditions
\begin{itemize}
  \item[a)] $\Phi$ est fermée dans l'ens. ordonné $\mathfrak{F}_{\mathrm{él}}$
  \item[b)] Si \struck{$F_0, F'_0 \in$} $\mathfrak{F}_1, \mathfrak{F}_2 \in \Phi$, \add{$\mathfrak{F}_1 \neq \mathfrak{F}_2$} alors $\mathfrak{F}_1$, $\mathfrak{F}_2$ sont compatibles, i.e. $\Phi$ est une \emph{configuration}
  \item[c)] $\Phi$ est « loc. fini » en un sens convenable.
\end{itemize}
\note{$\mathfrak{F}_1$, $\mathfrak{F}_2$ : il écrit ici deux éléments de $\mathfrak{F}_{\mathrm{él}}$ avec la même capitale cursive que $\mathfrak{F}$}

\ill{} si on veut exclure le cas des figures infinies, il faut aussi considérer dans $\mathfrak{F}_{\mathrm{él}}$ \ill{} donnée des parties loc. finies (i.e. dont la famille des supports dans $X$ est loc. finie) qui forment un antifiltre dans $\mathfrak{F}_{\mathrm{él}}$.

Inversement, tout $\Phi$ satisfaisant a) b) c) provient d'une figure et d'une seule.

\page{9}
Ainsi, on est arrivé à la structure suivante
\begin{itemize}
  \item[(1)] $\mathfrak{F}_{\mathrm{él}}$ ens. des « figures élémentaires »
  \item[(2)] relation d'ordre sur $\mathfrak{F}_{\mathrm{él}}$, $\leq$ (inclusion des figures)
  \item[(3)] relation symétrique réflexive $\mathcal{R}$ \struck{\ill{}} « $F$ et $G$ \emph{compatibles} ». \struck{\ill{}} \add{\emph{Configuration} (élémentaire) partie de $\mathfrak{F}_{\mathrm{él}}$ formée d'él. deux à deux compatibles}
  \item[(4)] \struck{Antifiltre sur $\mathfrak{F}_{\mathrm{él}}$ des parties} Parmi les configurations, \struck{\ill{}} \struck{\ill{}} celles qui sont « loc. finies », \struck{un} \add{\ill{} de l'ens.}
  \[
    \mathrm{Conf}_{\mathrm{l.f.}} \subset \mathfrak{P}(\mathfrak{F}_{\mathrm{él}})
  \]
\end{itemize}
\note{les nombres (1) à (4) sont cerclés sur la page. La définition de « configuration » dans (3) est écrite au-dessus et à droite de la ligne, par-dessus un mot biffé}
\marginal{\ill{} deux à deux compatibles \ill{}}
\note{note oblique dans la marge gauche, en face de (3) et (4), lue en partie}

\struck{Avec} propriétés d'antifiltre :
\begin{itemize}
  \item[a)] Toute configuration réduite à une seule figure él. est loc. finie
  \item[b)] Si $\Phi$ est une configuration, $\Phi$ loc. finie ssi la \struck{\ill{}} partie fermée (pour la rel. d'ordre $\leq$) \ill{}
  \item[c)] Si $\Phi$ est une configuration, \struck{$\Phi$, $\Psi$} loc. finie dans \uncertain{la}, \uncertain{toute} $\Psi \subset \Phi$ est loc. finie
  \item[d)] Si $\Phi$ est une configuration, $\Phi = \bigcup \Phi_i$ (réunion finie) où les $\Phi_i$ sont loc. finies, alors $\Phi$ \uncertain{aussi}.
\end{itemize}
\note{au-dessus de c), un mot (« l'\ill{} ») rattaché à la fin de b), qui porte peut-être sur la partie fermée engendrée par $\Phi$}
\marginal{(les \uncertain{configurations} \ill{} « \ill{} ») : loc. finies \ill{} \ill{} $X$ \ill{} finies}
\note{seconde note oblique dans la marge gauche, en face de a)--d), lue par fragments}

\page{10}
On définit dans l'ens. $\mathfrak{F}$ des \emph{figures}, comme l'ens. des parties $\Phi = \Sigma'_F$ de $\mathfrak{F}_{\mathrm{él}}$, qui sont
\begin{itemize}
  \item[a)] des configurations
  \item[b)] fermées pour la relation d'ordre
  \item[c)] loc. finies.
\end{itemize}
\note{un « $F$ » est écrit au-dessus de « figures »}
\emph{Figures finies} : parties finies de $\mathfrak{F}_{\mathrm{él}}$ satisfaisant a) et b).
\marginal{\uncertain{satisfaisant} a) \ill{} c) \ill{} \ill{} $X \ill{} \mathfrak{F}_{\mathrm{él}}$ \ill{} figures \ill{} \ill{}}
\note{deux notes obliques couvrent la marge gauche en face de ces lignes, en partie sur « \emph{finies} » ; elles ne sont lues que par fragments}

Notons encore une propriété, en plus des propriétés a) à d) \add{(p. 9)} ci-dessus, de la relation entre (2) et (3)
\note{(2) et (3) sont cerclés ; ils renvoient aux données (2), (3) de la page 9, dont viennent aussi les propriétés a) à d)}

\struck{Axiome d.}

\emph{Fig$_1$} Si $\mathbf{X}_1, \mathbf{X}_2 \leq \mathbf{X}$ sont des fig. él. alors $\{\mathbf{X}_1, \mathbf{X}_2\} \in \mathcal{R}$ i.e. $\{\mathbf{X}_1, \mathbf{X}_2\}$ est une configuration. [Dans l'ens. \add{des $\mathbf{X}'$} $\mathbf{X}' \leq \mathbf{X}$ est une configuration, et en vertu de a) b) c) donc est une \emph{figure} : la figure associée à $F$\ldots]
\note{à partir d'ici il note les figures élémentaires par un $X$ barré, rendu $\mathbf{X}$ ; le premier $\mathbf{X}_1$ est écrit sur un signe biffé. Le « $F$ » de la fin est sic, pour $\mathbf{X}$}
\marginal{\ill{} figure élémentaire associée \ill{}}
\note{note oblique dans la marge gauche, en face du crochet, lue en partie}

\page{11}
\emph{Fig$_2$} \add{Soit $\mathbf{X} \in \mathfrak{F}_{\mathrm{él}}$,} Considérons l'ens. ordonné des $\mathbf{X}' \leq \mathbf{X}$ est \emph{de dim. comb. finie}, i.e. il existe un entier $n$, tel que les parties finies \struck{\ill{}} majorées par $\mathbf{X}$, \add{et tot. ordonnées} sont de cardinal $\leq n$.
\note{l'ajout « Soit $\mathbf{X} \in \mathfrak{F}_{\mathrm{él}}$, » est écrit au-dessus de la ligne et renvoyé devant « Considérons » ; « et tot. ordonnées » est écrit au-dessus de la ligne. Fig$_1$ et Fig$_2$ sont réunis par un trait vertical dans la marge}

La \emph{dimension combinatoire} d'un $\mathbf{X} \in \mathfrak{F}_{\mathrm{él}}$ est le Sup des cardinaux des parties finies $\Phi$ précédentes, diminués de 1. Donc la dim combinatoire est $\geq 0$, et elle est nulle ssi $\mathbf{X}$ est un élément \emph{minimal} de $\mathfrak{F}_{\mathrm{él}}$.

Dans les exemples types, \uncertain{ce} les figures élémentaires minimales correspondent aux parties fermées de $X$ qui sont des \emph{variétés connexes}.

Si $\Phi \subset \mathfrak{F}_{\mathrm{él}}$ est une figure, sa \emph{dim combinatoire} est le dim combinatoire, en tant qu'ens. ordonné, i.e. c'est le sup des dim combin. $\mathbf{X}$, pour $\mathbf{X} \in \Phi$. Elle est nulle ssi $\Phi \neq \emptyset$ et \struck{$\Phi$ est formé de \ill{}} la relation

\page{12}
d'ordre sur $\Phi$ est discrète i.e. $\Phi \neq \emptyset$ et formé d'éléments minimaux de $\mathfrak{F}$.
\note{sic : $\mathfrak{F}$, pour $\mathfrak{F}_{\mathrm{él}}$}

Ce sont donc les \struck{ens.} \struck{configurations} parties de $\mathfrak{F}_{\mathrm{él}\,\mathrm{min}}$ qui ont les propriétés suivantes
\note{« formé de fig. él. » est écrit à droite, au bout de la ligne}
\begin{itemize}
  \item[a)] c'est une configuration, i.e. \add{mutuellement} \struck{dis} compatibles, et elle est non vide
  \item[b)] Elle est loc. finie.
\end{itemize}

Dans les exemples types, ceci correspond aux parties \add{fermées} $Y$ de $X$ qui sont des sous-variétés \add{non vides}, \uncertain{décomposées} en strates par les composantes connexes.

Ainsi, les figures de dim $\leq 0$ s'identifient les \emph{sous-variétés} (en sens strict, i.e. avec stratification triviale). La sous-variété vide correspond à la configuration vide.

\page{13}
\emph{Subdivision d'une figure}. \struck{\ill{}} On se reporte au cas-type. Soit une figure $F$, fermée de $X$ et de \struck{\ill{}} de ses strates $X_i$. Une \struck{subdivision} de figure $F'$ qui est une subdivision de $F$ est \struck{\ill{}} définie par les propriétés
\note{le $\mathfrak{F}$ en tête de la ligne suivante est barré ; « fermée de $X$ et de ses strates $X_i$ » : sic, la phrase n'est pas achevée}

a) $|F'| = |F|$ \qquad b) la relation d'équiv. dans $X = |F|$ définie par $F'$ est plus fine que celle définie par $F$.

La condition b) implique

b') $\forall\, X_i$, strate (fermée) de $X$ pour $F$, $X_i$ est réunion de strates de $X$ pour $F'$, et \add{elle} \struck{\ill{}} est $=$ équivalente. Ainsi, les strates pour $F'$ qui sont $\subset X_i$, forment une figure, qui \add{est une subdivision de $F_{X_i}$} \add{élémentaire} \struck{\ill{}} (la figure \add{induite} par $F$ sur $X_i$) Mais en fin de compte cette figure, de support $X_i$, n'est pas nécessairement élémentaire (une subdivision d'une figure élém. peut n'être pas élémentaire)
\note{« est une subdivision de $F_{X_i}$ », « élémentaire » et « induite » sont écrits au-dessus des lignes ; leur place exacte dans la phrase n'est pas sûre}

\page{14}
On va dire les choses autrement, en introduisant une autre relation d'ordre dans $\mathfrak{F}_{\mathrm{él}}$, plus fine que l'inclusion. \struck{\ill{}}

Soient $F$, $F'$ deux fig. \struck{élémentaires}. Je dis que $F' \ll F$ (que $F'$ \emph{raffine} $F$) si \struck{toute} \struck{cellule} strate \emph{ouverte} de $F'$ est contenue dans une \add{(unique, bien sûr)} strate ouverte de $F$. \emph{NB} On a aussi \struck{\ill{}} donc $|F'| \subset |F|$, mais pas nécessairement $|F'| = |F|$) C'est une relation \emph{d'ordre} sur l'ens. des figures, plus forte que celle d'inclusion (qui signifie \struck{\ill{}} que toute \add{strate} \struck{cellule} ouverte de $F'$ est une \add{strate} \struck{cellule} ouverte de $F$). Cette relation \ill{} \ill{} \struck{\ill{}} \ill{} \uncertain{se déduit} de la relation d'ordre induite sur l'ens. des figures élémentaires, $\mathfrak{F}_{\mathrm{él}}$, \ill{} on a
\note{dans la première formule, les indices et exposants de $F$, $F'$ sont surchargés (restes d'une notation $F_0$, $F'_0$ abandonnée) ; on garde la lecture finale $F' \ll F$. Au-dessus de la ligne, une parenthèse fermante renvoie après « $|F'| = |F|$ »}

$F'$ raffine $F$ ssi pour tout \struck{$F'_{X'}$} ($X' \in \Sigma_{F'}$), $\exists\, X \in \Sigma_F$, telle que $F'_{X'}$ raffine $F_X$

\page{15}
Ceci dit, une figure $F'$ est une \emph{subdivision} d'une figure $F$ ssi elle \emph{raffine} $F'$, et \add{si} elle a \emph{même support}. Mais (cf plus bas) cela équivaut : \add{ssi} $F$ est \emph{maximale} \add{par} \struck{parmi} les figures qui raffinent $F$. \struck{\ill{}} Ainsi, on trouve \ldots{} \ill{} \ill{} de structure :
\note{sic : « raffine $F'$ », pour « raffine $F$ ». Cette phrase est très corrigée : « ssi », « Mais (cf plus bas) », « par » sont écrits au-dessus de la ligne ; « la relation $\leq$, \struck{pour} tout \ill{} les » est intercalé et n'est lu qu'en partie}

\begin{itemize}
  \item[(5)] Une relation d'ordre $\ll$ sur $\mathfrak{F}_{\mathrm{él}}$ (« raffinement »), impliquée par $\leq$, d'où une relation, pour deux figures $F \subset \mathfrak{F}_{\mathrm{él}}$ et $F' \subset \mathfrak{F}_{\mathrm{él}}$, \emph{$F$ raffine $F'$} :
  \[
    \struck{\ill{}}\ \forall\, X \in F, \quad \exists\, X' \in F' \text{ telle que } X \ll X'.
  \]
\end{itemize}
\note{(5) est cerclé, à la suite des données (1) à (4) de la page 9. Devant $\forall$ et devant $F$, des signes sont raturés}

Ceci posé : $X \leq X'$ ssi \add{Pour $X, X' \in \mathfrak{F}_{\mathrm{él}}$, on a :} \struck{$F_{X'}$ \ill{}} $\{X, X'\}$ comp. et $X \ll X'$.

\struck{(6) On appellera subdivision d'une figure $F$, $F' \subset \mathfrak{F}_{\mathrm{él}}$, on a \ill{} des figures $F \ll F'$}

\struck{Conséquence $\Longrightarrow$ $F$, $F'$ compatibles, et}

\struck{$F \leq F'$}
\note{un (6) cerclé ouvrait ces lignes ; le passage est barré de plusieurs traits}
\marginal{Dans les structures (3) et (5), \ill{} structure (1) \ill{}}
\note{note oblique dans la marge gauche, en face du passage biffé ; (3), (5), (1) y sont cerclés}

Ceci posé, on dira qu'une figure $F'$ est une \emph{subdivision} d'une autre $F$, si c'est un \emph{raffinement}, et si elle est maximale, pour la relation $\subset$, parmi les raffinements de $F$.

\page{16}
\emph{Justification} (?) Revenons au cas type, \add{de $F$} $F'$ un raffinement, tel que $|F'| \neq |F|$. Montrons que $F'$ n'est pas \add{pour $\leq$,} maximal, parmi les raffinements de $F$. On va définir un raffinement \uncertain{canonique} \add{$F''$} de $F$, qui est \struck{\ill{}} une subdivision, et qui contient la figure $F'$. (donc \uncertain{strictement}, puisque $\underbrace{|F'|}_{X'} \neq \underbrace{|F|}_{X} = |F''|$) Les \struck{cellules} strates ouvertes de $F''$ seront
\begin{itemize}
  \item[a)] les strates \add{$U'_j$} ouvertes de $F'$
  \item[b)] Les comp. connexes \add{$V_{i\alpha}$} des \struck{\ill{}} $U_i \smallsetminus U_i \cap X'$ ($U_i$ est une strate ouverte de $F$) (i.e. chaque \uncertain{fois} que ce n'est pas $\emptyset$)
\end{itemize}
\note{« $U'_j$ » et « $V_{i\alpha}$ » sont écrits au-dessus des lignes a), b), rattachés par un trait ; « $U_i \cap X'$ » est souligné d'une accolade}
Il faut prouver les conditions \struck{\ill{}} de stratification, et pour commencer \struck{\ill{}} d'une
\[
  \partial V_{i\alpha} \overset{\mathrm{déf}}{=} \overline{V_{i\alpha}} \smallsetminus V_{i\alpha} \quad \text{est réunion de strates}
\]
Donc dire que $\partial V_{i\alpha} \cap U_i \subset$ \uncertain{Frontière} de $V_{i\alpha}$ dans $U_i$ \add{$\subset X' \cap U_i$} il y a un pb, qui a l'air intéressant, de topologie d'un espace stratifié \add{$(X' \cap U_i)$} dans une \emph{variété} (top. ou modérée) $(U_i)$
\note{la fin de la page est serrée et en partie surchargée ; « topologie » est écrit au-dessus d'un mot biffé}

\marginal{À \ill{} \ill{}, \ill{} \ill{} raffinement \ill{} \ill{} subdivision \ill{} \ill{} \ill{} $F'$ \ill{} \ill{} $X'$ \ill{} $U_i \cap F \cup$ \ill{} $F''$ \ill{} figure \ill{} $F'$ \ill{}}
\note{une longue note oblique occupe toute la marge gauche de la page ; elle n'est lue que par fragments}

\page{17}
Une \emph{contrée} est définie par la \add{donnée} \ill{}
\begin{itemize}
  \item[1°)] Un ensemble de « figures élémentaires », $\mathfrak{F}_{\mathrm{él}}$
  \item[2°)] Une relation d'ordre \struck{$\leq$} $\ll$ \struck{\ill{}} (raffinement) parmi les figures élémentaires
  \item[3°)] Une relation réflexive symétrique $\mathcal{R}$ (« compatibilité ») parmi \add{les} figures élémentaires
  \item[\uncertain{4°)}] \ill{} un antifiltre \uncertain{dans} l'ens. des parties de $\mathfrak{F}_{\mathrm{él}}$, appelées loc. « \uncertain{finies} » -- dans le cas de nos « contrées » finies, ce sont les parties \emph{finies} qui sont des « configurations » i.e. formées de figures él. deux à deux compatibles, satisfaisant à des conditions convenables, à dégager. Il y a déjà celle-ci :
\end{itemize}
\note{au-dessus de 1°) et 2°), des mots obliques (« \uncertain{régulière} » ?) et, entre les lignes 2°) et 3°), une insertion oblique « $\leq$ ou \ill{} » ne sont lus qu'en partie ; « compatibilité » est écrit sur un mot biffé ; le numéro « 4°) » est dans la note de marge et n'est pas sûr}
\marginal{\emph{NB} \ill{} \ill{} \ill{} \ill{} 4°) \ill{} \ill{} configuration \ill{} \ill{} \ill{} contrée finie \ill{} plutôt « configuration » \ill{} figure \ill{}}
\note{longue note oblique dans la marge gauche, lue par fragments}

\emph{Contr$_1$} Si \add{Soient} deux fig. élémentaires $\mathbf{X}$, $X'$, on \uncertain{pose} $X' \leq \mathbf{X}$ ssi (déf) $X' \ll X$, et $X'$ et $X$ sont comp. alors on trouve une relation \emph{d'ordre},

De plus, en termes de cette relation d'ordre, on a les propriétés Fig$_1$ \add{et $\mathrm{Fig}_2$}
\[
  (\Longleftrightarrow [\text{si } X', X'' \ll X, \text{ sont compatibles à } X \text{, ils sont compatibles entre eux}])
\]
Fig$_2$ (dim combinatoire finie) On va
\note{« et Fig$_2$ » (ou « Fig$_{1,2}$ ») est écrit sous « Fig$_1$ » ; « et à $X$, $X'$ » est écrit au-dessus de « sont compatibles ». Un double trait vertical marque la ligne entre crochets dans la marge}

\page{18}
la \struck{raffinement} \add{renforcer}. \struck{--}
\note{« renforcer » est écrit au-dessus du mot biffé ; il continue « On va » de la page 17}

\struck{Contr$_2$ Si $\mathbf{X} \in \mathfrak{F}_{\mathrm{él}}$, dans l'ens. \struck{\ill{}} \struck{\ill{}} $\{X' \in \mathfrak{F}_{\mathrm{él}} \mid X' \ll \mathfrak{F}_{\mathrm{él}}\}$}
\note{ces lignes sont barrées de traits obliques ; sic, « $X' \ll \mathfrak{F}_{\mathrm{él}}$ »}

\emph{Contr$_2$} Pour $\forall\, \mathbf{X} \in \mathfrak{F}_{\mathrm{él}}$, considérons la dim comb. (pour $\leq$) des $X' \in \mathfrak{F}_{\mathrm{él}}$ \struck{tel que} qui raffinent $\mathbf{X}$. Alors les dim. comb. de ces $X'$ est majorée.

\struck{\ill{}} Leur borne sup. s'appelle la \emph{dimension} \add{(géom.)} de $X$. La dim géom d'une figure, est par déf. le sup des \struck{\ill{}} dim géom. des figures élémentaires constituantes. C'est donc le sup \add{des $n$} pour \add{les} \uncertain{longueurs} des \uncertain{chaînes} $(X_0, X_1, \ldots, X_n)$ avec
\[
  X_0 < X_1 < \cdots < X_{n-1} \ll X_n \in F .
\]
\marginal{\struck{Donc la dimension \ill{}}}
\note{le « $\ll$ » est écrit plus gras que les « $<$ » ; au bout de la ligne, une insertion serrée (« Elle \ill{} \ill{} \ill{} $F$ \ill{} $F$ raffine $F$, élément \ill{} subdivision ») n'est lue qu'en partie}

\marginal{\emph{(11 juin)}}
\note{date de sa main, soulignée, dans la marge gauche en face de ce paragraphe : « (11 juin) », c'est-à-dire le 11 juin 1986, le lendemain du 10 juin de la page 1}
On aurait pu dire \ill{} « \ill{} » \ill{} un \ill{} la notion de parties \add{(\struck{\ill{}})} \emph{constructibles}, et \struck{\ill{}} admissions \add{\ill{} figures \ill{}} \struck{\ill{}} \struck{\ill{}} \uncertain{constructibles} \struck{\ill{}} comme \struck{\ill{}} \ill{} \ill{} \add{d'une figure} \uncertain{de} \uncertain{variétés} \add{à une figure} \ill{} donnée \uncertain{par} \ill{} \ill{}
\[
  \ll, \quad \text{compatibilité sur } \mathfrak{F}_{\mathrm{él}} \struck{\ill{}}.
\]
\ill{} sérieusement. Pour la \uncertain{peine} ultérieurement, j'ai envie de pousser \struck{\ill{}} un peu dans la direction $\ll$, compatibilité
\marginal{\uncertain{Aussi} \ill{} \ill{} \ill{} figures \ill{} \ill{} \ill{} \ill{}}
\note{le bas de la page est un palimpseste : plusieurs lignes d'insertions obliques, en partie biffées, se chevauchent au-dessus du texte principal et dans la marge gauche ; seul le fil principal est donné, par fragments}

\page{19}
\uncertain{j'en} arrive à la notion de lieu, et d'interprétation des \uncertain{domaines} en termes de lieux.

Mais \uncertain{en} terminologie, au lieu de « figure élémentaire », je préfère dorénavant « \emph{multistrate} », plus parlant. Donc une figure quelconque est un ens. de multistrates mutuellement compatibles, fermé pour $\leq$, et loc. fini.

\struck{\ill{}} On appelle \struck{multi} \emph{lieux} de la contrée, un multistrate qui est minimal pour $\ll$, i.e. qui n'a d'autre raffinement que lui-même. \struck{C'est donc} \struck{C'est donc} \struck{Dans les} (Dans les exemples-types, ils \uncertain{sont} \uncertain{réduits} à une \ill{} \ill{} de parties de $\mathbf{X}$, réduits à une seule partie, qui est ponctuelle -- donc les lieux sont les pts de $X$ --) Ce sont donc aussi des multistrates minimaux pour $\leq$, \struck{\ill{}} \add{i.e. \uncertain{faits} des} « variétés » \struck{\ill{} non vides} -- ce sont les sous-variétés $0$-connexes qui n'ont pas de raffinement autre qu'elle-même.
\note{à droite des deux lignes biffées « C'est donc », quelques mots raturés dans la marge. $\mathbf{X}$ est ici l'$X$ barré ; « $0$-connexes » : lecture du chiffre sûre, du sens non}

Soit $\mathcal{L}$ l'ensemble des lieux,
\[
  \struck{\text{Pour tout}}\ \mathcal{L} \subset \mathcal{M} \quad (= \mathfrak{F}_{\mathrm{él}} = \text{ensemble des multistrates})
\]
\note{sous « Pour tout » biffé, un mot (« donc » ?) ; $\mathcal{M}$ est sa capitale cursive M}

\page{20}
\struck{toute} Pour tout multistrate $X$, soit $\mathcal{L}_X$ \add{$\in \mathfrak{P}(\mathcal{L})$} \struck{l'ens.} des lieux qui raffinent $X$. \marginal{$\widehat{\mathcal{L}}_X$} et soit
\[
  \widehat{\mathcal{L}}_X \subset \mathfrak{P}(X) \qquad \text{i.e.} \qquad \widehat{\mathcal{L}}_X \in \mathfrak{P}(\mathfrak{P}(X))
\]
l'ensemble des \struck{\ill{}} $\mathcal{L}_Y$ pour $Y \leq X$ ($Y$ sous-multistrate de $X$). \struck{Considérons} \struck{les} J'aimerais reconstituer les structures $\ll$ et comp. sur $\mathcal{M}$, en termes de l'application
\[
  X \longmapsto \widehat{\mathcal{L}}_X \qquad \mathcal{M} \longrightarrow \mathfrak{P}(\mathfrak{P}(X)) .
\]
\note{le chapeau de $\widehat{\mathcal{L}}_X$ est écrit par-dessus un signe raturé ; $\mathfrak{P}(X)$ : sic, on attendrait $\mathfrak{P}(\mathcal{L})$, comme plus bas}

\struck{\emph{Contr$_3$}} Notons que pour $\mathcal{F} \subset \mathfrak{P}(\mathcal{M})$, \uncertain{sous-ens.} figures, on pose
\[
  \struck{\ill{}}\ \mathcal{L}_{\mathcal{F}} = \text{ens. des lieux qui raffinent } \mathcal{F} \quad \text{i.e.} \quad \mathcal{L}_{\mathcal{F}} = \bigcup_{X \in \mathcal{F}} \mathcal{L}_X \in \mathcal{L}
\]
\[
  \widehat{\mathcal{L}}_{\mathcal{F}} = \{\mathcal{L}_X \mid X \in \mathcal{F}\} \in \mathfrak{P}(\mathcal{L}) \qquad \text{i.e.} \in \mathfrak{P}(\mathfrak{P}(\mathcal{L}))
\]
qui nous donne
\[
  \mathcal{F} \longmapsto \widehat{\mathcal{L}}_{\mathcal{F}} \qquad \mathcal{F} \longrightarrow \mathfrak{P}(\mathfrak{P}(X))
\]
\note{il note ici une figure $\mathcal{F}$ par sa capitale cursive F ; « sous-ens. » est écrit au-dessus de la ligne, lecture douteuse. Sic : « $\in \mathcal{L}$ » pour $\subset \mathcal{L}$, « $\in \mathfrak{P}(\mathcal{L})$ » corrigé en « $\in \mathfrak{P}(\mathfrak{P}(\mathcal{L}))$ », et $\mathfrak{P}(\mathfrak{P}(X))$ en fin de page}

\struck{\emph{Contr$_3$} a) Pour que $X \ll Y$ \struck{faut-il} $X, Y \in$}

\struck{Soient $X$, $Y$ deux multistrates. \struck{Si} Supposons $X \ll Y$. Alors pour tout $X_i \leq X$, \ill{}}
\note{ces dernières lignes sont barrées de traits obliques ; la page s'arrête là}

\end{document}
